\documentclass[12pt]{beamer}
\usepackage{../ceai-beamer}
% \usepackage[handoutnotes,letterpaper]{../ceai-beamer}

\title{CE 488/5588: Applied Machine Learning}
\subtitle{Topic 2: Data Types and Revisiting Linear Algebra}
\author{}
\date{}

\begin{document}

\begin{frame}
  \titlepage
\end{frame}

\begin{frame}{Outline}
  \tableofcontents
\end{frame}

% ==========================================================
% Observations as Data
% ==========================================================
\section{Observations as Data}

\begin{frame}{From observation to induction}
  \begin{conceptbox}{Why data matters}
    Inductive reasoning depends on observations. For humans, sensing and cognition are intertwined.
    For machines, observation becomes explicit data that must be stored, interpreted, and acted on.
  \end{conceptbox}
  \begin{itemize}
    \item Learning systems extract patterns from observation to build general laws.
    \item Data quality and representation set the ceiling for model capability.
  \end{itemize}
\end{frame}

\begin{frame}{Three subsystems of an intelligent machine}
  \begin{itemize}
    \item Perception: sensors capture the environment.
    \item Computing: the system processes and interprets sensed data.
    \item Actuation: decisions are translated into actions.
  \end{itemize}
  % \vspace{0.6em}
  % \centering
  % \LARGE \textbf{[Figure Placeholder]}
  % \vspace{0.8em}
  % \normalsize
  % \textit{Perception--computation--actuation pipeline with feedback loop.}
\end{frame}

\begin{frame}{From sensing to data}
  \begin{itemize}
    \item Sensors produce raw data, typically numeric and high dimensional.
    \item Data can be large-scale in size, complexity, and dimensionality.
    \item Data representation determines what learning is feasible.
  \end{itemize}
\end{frame}

\begin{frame}{Data vs. information}
  \begin{conceptbox}{Operational distinction}
    Data are raw measurements; information is data with context or processing applied.
  \end{conceptbox}
  \begin{itemize}
    \item Many decisions are categorical: yes/no, class labels, or ranked levels.
    \item The ML pipeline must preserve the semantics of the data type.
  \end{itemize}
\end{frame}

\begin{frame}{Dimensionality preview}
  \begin{itemize}
    \item Data can be low-dimensional (scalar values) or high-dimensional (vectors, matrices, tensors).
    \item Dimension affects storage, visualization, and model design.
    \item We will revisit dimensionality using a concrete-strength monitoring example.
  \end{itemize}
\end{frame}

% ==========================================================
% Types of Data
% ==========================================================
\section{Types of Data}

\begin{frame}{Two major data types}
  \begin{itemize}
    \item Categorical data: categories such as class labels, material types, or damage states.
      \begin{itemize}
        \item Can appear as outputs ($y$) or as inputs (features in $\bm{x}$).
      \end{itemize}
    \item Numeric data: values on numerical scales (interval or ratio).
      \begin{itemize}
        \item Often require scaling (standardization/normalization) before training.
      \end{itemize}
  \end{itemize}
  \vspace{0.3em}
  \textit{Measurement scales provide a formal framework for these categories.}
  \bhref{https://en.wikipedia.org/wiki/Level_of_measurement}{Measurement scales}
\end{frame}

\begin{frame}{Categorical variables are not only labels}
  \begin{itemize}
    \item In supervised learning, we often write a data set as $\mathcal{D}=\{(\bm{x}_i,y_i)\}_{i=1}^n$.
    \item A categorical variable can be:
      \begin{itemize}
        \item a \textbf{label} $y_i$ (e.g., damage state), or
        \item an \textbf{input feature} (e.g., material type, sensor ID, site category).
      \end{itemize}
    \item Common encodings: one-hot for nominal categories; integer encoding for ordinal categories.
  \end{itemize}
\end{frame}

\begin{frame}{Categorical data: nominal}
  \begin{conceptbox}{Nominal data}
    Labels that classify objects into mutually exclusive categories, with no inherent order.
  \end{conceptbox}
  \begin{itemize}
    \item Fruit types: apple, pear, orange.
    \item Materials: concrete, steel, timber.
    \item Planets: Earth, Mars, Venus.
  \end{itemize}
\end{frame}

\begin{frame}{Categorical data: ordinal}
  \begin{conceptbox}{Ordinal data}
    Categories with an inherent order or ranking, but no meaningful distance between ranks.
  \end{conceptbox}
  \begin{itemize}
    \item Pollution: none, light, severe, lethal.
    \item Damage: intact, minor, major, collapse.
  \end{itemize}
\end{frame}

\begin{frame}{Numeric data: interval}
  \begin{conceptbox}{Interval data}
    Differences are meaningful and equal, but the zero point is arbitrary.
  \end{conceptbox}
  \begin{itemize}
    \item Temperature in Celsius or Fahrenheit.
    \item Velocity or acceleration expressed without physical zero reference.
  \end{itemize}
\end{frame}

\begin{frame}{Numeric data: ratio}
  \begin{conceptbox}{Ratio data}
    Values have a physical zero, making ratios meaningful.
  \end{conceptbox}
  \begin{itemize}
    \item Temperature in Kelvin (absolute zero as origin).
    \item Mass, length, duration, energy, electric charge.
  \end{itemize}
  \bhref{https://www.nist.gov/si-redefinition/kelvin-introduction}{NIST Kelvin}
\end{frame}

\begin{frame}{Discrete vs. continuous}
  \begin{itemize}
    \item Discrete data are countable (integers in $\mathcal{N}$).
    \item Continuous data vary over a real interval (values in $\mathcal{R}$).
    \item Most continuous signals are digitized for processing.
  \end{itemize}
\end{frame}

\begin{frame}{Levels of measurement}
  \begin{figure}
    \centering
    \includegraphics[width=0.78\textwidth]{Figure0102.jpg}
    \caption{Levels of measurement.}
  \end{figure}
\end{frame}

\begin{frame}{Temperature as interval and ratio}
  \begin{examplebox}{Interval vs. ratio temperature}
    Celsius and Fahrenheit have arbitrary zero points, so ratios are meaningless.
    Kelvin defines absolute zero, so ratios are physically meaningful.
  \end{examplebox}
  \begin{itemize}
    \item $40\,^{\circ}F$ is not twice as hot as $20\,^{\circ}F$.
    \item $40\,K$ is twice $20\,K$ due to absolute zero.
  \end{itemize}
\end{frame}

\begin{frame}{Transforming scales}
  \begin{itemize}
    \item Physical quantities can be mapped between interval and ratio scales.
    \item Example: $y = \log(x)$ maps a positive ratio value $x>0$ to an interval-like scale.
    \item Example: $x = \exp(y)$ maps an interval value back to a positive ratio scale.
    \item These transforms are often used as \textit{feature engineering} (to reduce skew, stabilize variance).
  \end{itemize}
\end{frame}

\begin{frame}{A translation layer: narrative $\rightarrow$ ML notation}
  \begin{itemize}
    \item We model a data set as $n$ samples: $\mathcal{D}=\{(\bm{x}_i,y_i)\}_{i=1}^n$.
    \item $\bm{x}_i \in \mathbb{R}^d$ is the feature vector (inputs); $y_i$ is the target (output).
    \item Tabular inputs are often stored in a matrix $\bm{X}\in \mathbb{R}^{n\times d}$, where each row is one sample.
  \end{itemize}
\end{frame}

\begin{frame}{Data $\rightarrow$ features (preprocessing)}
  \begin{itemize}
    \item Raw measurements may need cleaning and scaling before they become features.
    \item Standardization (z-score): $z=\dfrac{x-\mu}{\sigma}$ (approximately zero mean, unit variance).
    \item Normalization (min--max): $x' = \dfrac{x-x_{\min}}{x_{\max}-x_{\min}}$ (rescale to $[0,1]$).
  \end{itemize}
\end{frame}

% ==========================================================
% Notations and Dimensionality
% ==========================================================
\section{Notations and Dimensionality}

\begin{frame}{Digital vs. analog sensing}
  \begin{itemize}
    \item Analog sensors output continuous values.
    \item Digital sensors output discrete samples.
    \item Most signals are digitized before ML processing.
  \end{itemize}
\end{frame}

\begin{frame}{Thought experiment: concrete strength}
  \begin{examplebox}{Monitoring a hydraulic dam}
    Concrete strength is measured at day 28 and monitored across space and time.
    The experiment explores how data dimensionality grows with locations and time.
  \end{examplebox}
  \begin{itemize}
    \item Each measurement at a fixed location and time is a scalar.
    \item Complexity arises when we aggregate many measurements into one sample.
  \end{itemize}
\end{frame}

\begin{frame}{Measurement vs. sample dimension vs. data set size}
  \begin{itemize}
    \item \textbf{Scalar measurement:} one number at one time and one location.
    \item \textbf{Sample/feature dimension $d$:} how many numbers we stack into one input $\bm{x}_i \in \mathbb{R}^d$.
      \begin{itemize}
        \item Example: a $10\times 10$ grid can be flattened into $d=100$ features.
      \end{itemize}
    \item \textbf{Data set size $n$:} how many samples we have (times, locations, experiments).
  \end{itemize}
\end{frame}

\begin{frame}{Measurement schemes (S1--S3)}
  \scriptsize
  \renewcommand{\arraystretch}{1.2}
  \setlength{\tabcolsep}{3pt}
  \begin{tabular}{|p{1.2cm}|p{2.2cm}|p{2.2cm}|p{3.cm}|}
    \hline
    Scheme & Description & Data size & Interpretation \\
    \hline
    S1 & Single measurement at one location, day 28 & $1$ & Single scalar value \\ \hline
    S2 & Ten measurements at one location, day 28 & $1 \times 10$ & 1D sample set \\ \hline
    S3 & One location over 25 months & $1 \times 25$ & Time series of scalars \\
    \hline
  \end{tabular}
\end{frame}

\begin{frame}{Measurement schemes (S4--S5)}
  \scriptsize
  \renewcommand{\arraystretch}{1.2}
  \setlength{\tabcolsep}{3pt}
  \begin{tabular}{|p{1.2cm}|p{3.2cm}|p{2.0cm}|p{3.cm}|}
    \hline
    Scheme & Description & Data size & Interpretation \\
    \hline
    S4 & $10 \times 10$ grid at day 28 & $10 \times 10$ & Matrix (image-like) \\ \hline
    S5 & $10 \times 10$ grid over 25 months & $10 \times 10 \times 25$ & Time series of matrices \\
    \hline
  \end{tabular}
\end{frame}

\begin{frame}{Vector notation for S2 and S3}
  \begin{itemize}
    \item S2 and S3 can be represented as vectors or time series.
    \item Vector length equals the number of observations, but each data point is still 1D.
    \item Time series visualization requires a time axis in 2D space.
  \end{itemize}
\end{frame}

\begin{frame}{Matrix notation for spatial grids}
  \begin{itemize}
    \item S4 measurements form a matrix $\bm{F}$ of size $10 \times 10$.
    \item Each entry is a scalar value at a location.
  \end{itemize}
  \[
  \bm{F} =
  \begin{bmatrix}
  4.2 & 4.1 & \cdots & 3.9 \\
  \vdots & \vdots & \ddots & \vdots \\
  4.3 & 3.0 & \cdots & 4.1
  \end{bmatrix}
  \]
\end{frame}

\begin{frame}{S4 visualization}
  \begin{itemize}
    \item A matrix can be visualized as a surface plot in 3D.
    \item Axes: $x$ and $y$ locations, vertical axis is strength $f'_c$.
  \end{itemize}
  \vspace{0.5em}
  % \centering
  % \LARGE \textbf{[Figure Placeholder]}
  % \vspace{0.8em}
  % \normalsize
  % \textit{Surface plot of a $10 \times 10$ grid of strength values.}
\end{frame}

\begin{frame}{Nested matrices for S5}
  \begin{itemize}
    \item S5 produces a time series of matrices.
    \item A nested set of matrices represents measurement snapshots over time.
  \end{itemize}
  \[
  \bm{F} = \left\{\begin{bmatrix}
  4.2 & 4.1 & \cdots & 3.9 \\
  \vdots & \vdots & \ddots & \vdots \\
  4.3 & 3.0 & \cdots & 4.1
  \end{bmatrix}, \ldots, \begin{bmatrix}
  4.0 & 4.3 & \cdots & 4.4 \\
  \vdots & \vdots & \ddots & \vdots \\
  3.3 & 4.0 & \cdots & 3.9
  \end{bmatrix}\right\}
  \]
\end{frame}

\begin{frame}{Visualization challenge}
  \begin{itemize}
    \item Time-series matrices require a fourth axis for visualization.
    \item Plotting stacked surfaces becomes difficult to interpret.
    \item This motivates compact, high-dimensional representations.
  \end{itemize}
\end{frame}

\begin{frame}{Vector-valued measurements}
  \begin{examplebox}{Multiple tests per location}
    Concrete strength is measured using three techniques, yielding a 3D vector per location.
  \end{examplebox}
  \[
  \bm{f} = \{ (4.1, 3.9, 4.0)^T, (2.0, 2.8, 3.0)^T, \ldots \}
  \]
\end{frame}

\begin{frame}{Vector data as a matrix}
  \begin{itemize}
    \item Each column is a 3D vector (one location).
    \item Matrix rows represent measurement modalities.
  \end{itemize}
  \[
  \bm{F} =
  \begin{bmatrix}
  4.1 & 2.0 & \cdots & 4.1 \\
  3.9 & 2.8 & \cdots & 5.0 \\
  4.0 & 3.0 & \cdots & 4.2
  \end{bmatrix}
  \]
\end{frame}

\begin{frame}{Vector data over time}
  \begin{itemize}
    \item Monitoring over time yields a set of matrices.
    \item The structure is similar to S5, but the semantics differ.
    \item Each column is a data point with multiple attributes.
  \end{itemize}
\end{frame}

\begin{frame}{Toward tensors}
  \begin{itemize}
    \item High-dimensional data exceed matrices and vectors.
    \item Tensors are the natural representation for multi-axis data (space $\times$ time $\times$ channels).
    \item We revisit tensors (definition, notation, examples) in the \textbf{Tensors in Machine Learning} section.
  \end{itemize}
\end{frame}

% ==========================================================
% Imagery Data
% ==========================================================
\section{Special Data: Imagery Data}

\begin{frame}{Why imagery matters}
  \begin{itemize}
    \item Visual assessment dominates structural inspection and damage evaluation.
    \item Computer vision enables automated inspection with robotics.
    \item Remote sensing delivers large-scale imagery from satellites.
  \end{itemize}
  \bhref{https://www.earthdata.nasa.gov/learn/backgrounders/remote-sensing}{NASA Remote Sensing}
\end{frame}

\begin{frame}{Image storage basics}
  \begin{conceptbox}{Image as a matrix}
    An image is stored as a matrix of pixel values. A $100 \times 200$ image has
    100 rows and 200 columns of pixels.
  \end{conceptbox}
  \begin{itemize}
    \item Pixel = data point at a location in the imaging plane.
    \item Matrix size determines spatial resolution.
  \end{itemize}
\end{frame}

\begin{frame}{Practical tensor shapes for images and videos}
Define an image as $I$, and $H, W, C, T$ as its image height, width, thickness, and time-duration. 
  \begin{itemize}
    \item Grayscale image: $I_{(H,W)}$.
    \item Color image (RGB): $I_{(H,W,C)}$ with $C=3$.
    \item Video (time series of frames): $I_{(T,H,W,C)}$.
    \item Mini-batch adds a leading batch dimension $N$: $I_{(N,H,W,C)}$ or $I_{(N,T,H,W,C)}$.
    \item Terminology note: tensor \textit{order} (\# axes) is different from \textit{matrix rank} in linear algebra.
  \end{itemize}
\end{frame}

\begin{frame}{Binary image representation}
  \begin{itemize}
    \item Pixel values are binary: 0 (black) or 1 (white).
    \item Stored as a 2D matrix of zeros and ones.
  \end{itemize}
  \begin{figure}
    \centering
    \includegraphics[width=0.35\textwidth]{Figure0104a.png}
    \caption{Binary image example.}
  \end{figure}
\end{frame}

\begin{frame}{Grayscale image representation}
  \begin{itemize}
    \item Pixel values range continuously (e.g., 0 to 1).
    \item Often stored as integers 0--255 for 256 gray levels.
  \end{itemize}
  \begin{figure}
    \centering
    \includegraphics[width=0.35\textwidth]{Figure0104b.png}
    \caption{Grayscale image example.}
  \end{figure}
\end{frame}

\begin{frame}{Color image representation}
  \begin{itemize}
    \item Color images store three channels: red, green, blue.
    \item Each channel is a 2D matrix; together they form a 3D tensor.
  \end{itemize}
  \begin{figure}
    \centering
    \includegraphics[width=0.35\textwidth]{Figure0104c.png}
    \caption{Color image example.}
  \end{figure}
\end{frame}

\begin{frame}{Imagery and structural sensing}
  \begin{itemize}
    \item A grid of scalar measurements resembles a grayscale image.
    \item Multiple sensing modalities resemble RGB channels.
    \item Imagery offers intuition for matrix and tensor data structures.
  \end{itemize}
\end{frame}

\begin{frame}{Video data}
  \begin{itemize}
    \item Video is a time series of image frames.
    \item Adds a temporal dimension to image tensors.
  \end{itemize}
  \vspace{0.5em}
  \centering
  % \LARGE \textbf{[Figure Placeholder]}
  % \vspace{0.8em}
  % \normalsize
  % \textit{Stacked image frames along a time axis to form a 4D tensor.}
\end{frame}

% ==========================================================
% Linear Algebra Refresher
% ==========================================================
\section{Basic Linear Algebra Notations and Concepts}

\begin{frame}{Vectors in $\mathbb{R}^n$}
  \begin{itemize}
    \item A vector is an ordered list of scalars.
    \item Column vector notation is standard for ML.
  \end{itemize}
  \[
  \bm{v} =
  \begin{pmatrix}
  v_1 \\
  v_2 \\
  \vdots \\
  v_n
  \end{pmatrix}
  \]
\end{frame}

\begin{frame}{Matrices}
  \begin{itemize}
    \item A matrix is a 2D array of scalars.
    \item Size is $m \times n$ (rows by columns).
  \end{itemize}
  \[
  \bm{A} =
  \begin{pmatrix}
  a_{11} & a_{12} & \cdots & a_{1n} \\
  a_{21} & a_{22} & \cdots & a_{2n} \\
  \vdots & \vdots & \ddots & \vdots \\
  a_{m1} & a_{m2} & \cdots & a_{mn}
  \end{pmatrix}
  \]
\end{frame}

\begin{frame}{Symmetric matrices}
  \begin{conceptbox}{Definition}
    A matrix $\bm{A}$ is symmetric if $\bm{A} = \bm{A}^\top$.
  \end{conceptbox}
  \begin{itemize}
    \item Symmetry implies $a_{ij} = a_{ji}$.
    \item Common in covariance matrices and quadratic forms.
  \end{itemize}
\end{frame}

\begin{frame}{Identity matrices}
  \begin{itemize}
    \item The identity matrix $\bm{I}$ leaves vectors unchanged.
    \item It is the multiplicative identity for matrices.
  \end{itemize}
  \[
  \bm{I} =
  \begin{pmatrix}
  1 & 0 & \cdots & 0 \\
  0 & 1 & \cdots & 0 \\
  \vdots & \vdots & \ddots & \vdots \\
  0 & 0 & \cdots & 1
  \end{pmatrix}
  \]
\end{frame}

\begin{frame}{Sparse matrices}
  \begin{itemize}
    \item Sparse matrices contain mostly zeros.
    \item Storage and computation can be optimized using sparsity.
  \end{itemize}
  \[
  \bm{S} =
  \begin{pmatrix}
  0 & 0 & 1 \\
  0 & 0 & 0 \\
  2 & 0 & 0
  \end{pmatrix}
  \]
\end{frame}

\begin{frame}{Rank of a matrix}
  \begin{itemize}
    \item Rank is the dimension of the column space (or row space).
    \item Full rank implies maximum independent columns.
  \end{itemize}
\end{frame}

\begin{frame}{Determinant and invertibility}
  \begin{itemize}
    \item The determinant summarizes scaling and invertibility.
    \item A matrix is invertible if $\det(\bm{A}) \ne 0$.
  \end{itemize}
  \[
  \det(\bm{A}) \ne 0 \quad \Longleftrightarrow \quad \bm{A}^{-1} \text{ exists}
  \]
\end{frame}

\begin{frame}{Matrix inverse}
  \begin{itemize}
    \item The inverse satisfies $\bm{A} \bm{A}^{-1} = \bm{I}$.
    \item Used to solve linear systems when $\bm{A}$ is square and nonsingular.
  \end{itemize}
  \[
  \bm{A} \bm{A}^{-1} = \bm{A}^{-1} \bm{A} = \bm{I}
  \]
\end{frame}

\begin{frame}{Linear systems}
  \begin{itemize}
    \item Linear system: $\bm{A} \bm{x} = \bm{b}$.
    \item Unique solution if $\det(\bm{A}) \ne 0$.
  \end{itemize}
\end{frame}

\begin{frame}{Conditions for unique solutions}
  \[
  \det(\bm{A}) \ne 0 \Longleftrightarrow \mathrm{rank}(A) = n
  \Longleftrightarrow \text{unique solution}
  \]
  \begin{itemize}
    \item If $\det(\bm{A}) = 0$, the system may have no or infinite solutions.
  \end{itemize}
\end{frame}

\begin{frame}{Rectangular matrices}
  \begin{itemize}
    \item Many ML problems use $m \times n$ matrices with $m \ne n$.
    \item Exact inverses do not exist; pseudo-inverses are used instead.
  \end{itemize}
\end{frame}

\begin{frame}{Pseudo left inverse}
  \begin{itemize}
    \item If $m \ge n$ and $\mathrm{rank}(A) = n$:
  \end{itemize}
  \[
  A_\ell^+ = (A^\top A)^{-1} A^\top, \quad A_\ell^+ A = I_n
  \]
\end{frame}

\begin{frame}{Pseudo right inverse}
  \begin{itemize}
    \item If $m \le n$ and $\mathrm{rank}(A) = m$:
  \end{itemize}
  \[
  A_r^+ = A^\top (A A^\top)^{-1}, \quad A A_r^+ = I_m
  \]
\end{frame}

\begin{frame}{Moore-Penrose pseudo-inverse}
  \begin{itemize}
    \item Generalized inverse for arbitrary matrices.
    \item Produces minimum-norm least-squares solutions.
  \end{itemize}
  \[
  \bm{x} = A^+ \bm{b}
  \]
  \bhref{https://en.wikipedia.org/wiki/Moore\%E2\%80\%93Penrose_inverse}{Moore-Penrose inverse}
\end{frame}

\begin{frame}{Linear regression setup}
  \begin{itemize}
    \item Regression solves $X \boldsymbol{\beta} = \mathbf{y}$.
    \item $X$ has predictors, $\boldsymbol{\beta}$ are coefficients.
  \end{itemize}
\end{frame}

\begin{frame}{Least-squares solution}
  \begin{itemize}
    \item If $m \ge n$ and $X$ has full column rank:
  \end{itemize}
  \[
  \boldsymbol{\beta} = (X^\top X)^{-1} X^\top \mathbf{y}
  \]
\end{frame}

\begin{frame}{Underdetermined systems}
  \begin{itemize}
    \item If $m \le n$ and $X$ has full row rank:
  \end{itemize}
  \[
  \boldsymbol{\beta} = X^\top (X X^\top)^{-1} \mathbf{y}
  \]
\end{frame}

\begin{frame}{Unified least squares}
  \begin{itemize}
    \item Moore-Penrose pseudo-inverse unifies all cases.
  \end{itemize}
  \[
  \boldsymbol{\beta} = X^+ \mathbf{y}
  \]
\end{frame}

% ==========================================================
% Tensors
% ==========================================================
\section{Tensors in Machine Learning}

\begin{frame}{Why tensors}
  \begin{itemize}
    \item Tensors generalize scalars, vectors, and matrices.
    \item They represent images, videos, sensor grids, and model parameters.
  \end{itemize}
  \bhref{https://www.tensorflow.org/api_docs/python/tf/Tensor}{TensorFlow Tensor}
\end{frame}

\begin{frame}{Tensor definition}
  \begin{conceptbox}{Definition}
    A tensor is a multi-dimensional array of numerical values.
  \end{conceptbox}
  \begin{itemize}
    \item Rank (order) is the number of dimensions.
    \item In ML, tensors serve as data containers.
  \end{itemize}
\end{frame}

\begin{frame}{Tensor notation}
  \begin{itemize}
    \item Boldface for matrices ($\bm{T}$) and calligraphic for higher rank ($\mathcal{T}$).
    \item Components written as $T_{ijk\dots}$.
  \end{itemize}
  \[
  C_{ik} = A_{ij} B_{jk} \equiv \sum_j A_{ij} B_{jk}
  \]
\end{frame}

\begin{frame}{Examples by rank}
  \begin{itemize}
    \item Rank 0: scalar $T = 5.2$.
    \item Rank 1: vector $\bm{v} = [v_1, \ldots, v_n]$.
    \item Rank 2: matrix $\bm{W} = [W_{ij}]$.
  \end{itemize}
\end{frame}

\begin{frame}{Examples by rank (continued)}
  \begin{itemize}
    \item Rank 3: RGB image $\bm{I}_{ijk}$.
    \item Rank 4: video batch $\mathcal{B}_{ijkl}$.
  \end{itemize}
\end{frame}

\begin{frame}{Sensor data as tensors}
  \scriptsize
  \renewcommand{\arraystretch}{1.2}
  \setlength{\tabcolsep}{4pt}
  \centering
  \begin{tabular}{|p{2.0cm}|p{3.0cm}|p{4.cm}|}
    \hline
    Dimension & Structure & Examples \\
    \hline
    1D & $\{u_i\}$ & Temperature, stock prices \\
    2D & $\{u(x_i, y_j)\}$ & Grayscale images, grids \\
    3D & $\{u(x_i, y_j, t_k)\}$ & Video, MRI scans \\
    4D & $\{u(x_i, y_j, s_k, t_l)\}$ & Hyperspectral video \\
    \hline
  \end{tabular}
\end{frame}

\begin{frame}{Tensor applications in ML}
  \begin{itemize}
    \item Neural network weights and activations are tensors.
    \item GPUs operate efficiently on tensor operations.
    \item Autograd computes gradients like $\nabla_{\bm{W}} \mathcal{L}$.
  \end{itemize}
\end{frame}

% ==========================================================
% Closing
% ==========================================================
\section{Summary}

\begin{frame}{Key takeaways}
  \begin{itemize}
    \item Data types encode the semantics of measurements.
    \item Dimensionality shapes representation and visualization.
    \item Linear algebra tools support ML models and optimization.
    \item Tensors generalize data structures for modern ML.
  \end{itemize}
\end{frame}

\begin{frame}{Next steps}
  \begin{itemize}
    \item Apply data types to supervised and unsupervised learning tasks.
    \item Use linear algebra and tensors for model formulation.
  \end{itemize}
\end{frame}

\begin{frame}{Q\&A}
  \centering
  Questions?
\end{frame}

\end{document}
